\section{Types}
\label{sec:types}

This section assigns static meaning to the syntax  given in
\cref{fig:syntax}. The type system has three layers. First, the predicates 
$\Unr$, $\Mobile$, and $\Bounded$ classify the values, contexts, and sessions
that may be duplicated, transported across processes, or split; an auxiliary
predicate $\SNew$ restricts the sessions that may be created fresh. Second,
contexts are quotiented by equations that respect this classification, session
types are quotiented by a standard equational theory, and weakening is
controlled by a subcontext preorder. Third, the typing judgments for constants,
expressions, and processes use these classifications to enforce protocol
fidelity and ownership transfer.

Typing contexts combine assumptions with two constructors.
% \footnote{Inspired by
%   the tree-shaped contexts of \citet{DBLP:journals/bsl/OHearnP99}.}.
The ordered product
$\Ctx[1]\CtxSeq\Ctx[2]$ records that the resources in $\Ctx[1]$ must be used
before the resources in $\Ctx[2]$. The unordered product
$\Ctx[1]\CtxPar\Ctx[2]$ records independent resources. When a rule is
parametric in either constructor, we write $\ctxcolor{\Ctx[1]\CtxPatOp\Ctx[2]}$
for either $\Ctx[1]\CtxSeq\Ctx[2]$ or $\Ctx[1]\CtxPar\Ctx[2]$.

\begin{figure}
  \begin{mathpar}
    \RuleTeTerm \and
    \RuleTeWait \and
    \RuleTeDrop \and
    \RuleTeSeqOne \\
    \RuleTeRec \and
    \RuleTeBranch \and
    \RuleTeChoice \and
    \RuleTeSeqTwo \\
    \RuleTeBase \and
    \RuleTeArr \and
    \RuleTeAcq \and
    \RuleTeProd \and
    \RuleTeVariant \\
    \RuleTeCtxEmpty \and
    \RuleTeCtxVar \and
    \RuleTeCtxSeq \and
    \RuleTeCtxPar
  \end{mathpar}
  \caption{Mobile types and contexts ($\Mobile\T$, $\Mobile\Ctx$);
    bounded session types ($\Bounded\S$)}
  \label{fig:bounded-session-types}
\end{figure}

\Cref{fig:bounded-session-types} defines the predicates that classify sessions
and values for cross-process movement, together with the boundedness predicate
they rely on. A bounded session $\Bounded\S$ is a finite communication
obligation: a non-recursive bounded session terminates in $\STerm$, $\SWait$,
or $\SDrop$, while a recursive bounded session requires all finite unfoldings to
reach such a boundary consistently. A trailing session that is equivalent to
$\SSkip$ carries no communication obligation and is discharged when checking
boundedness (rule \RuleName{Te-Seq1}).

A mobile session has the form $\SSeq\SAcq\S$ with $\Bounded\S$. The leading
acquire marks the suffix of a remotely borrowed endpoint, and boundedness ensures that the
borrowed prefix can be completed without consuming an unbounded residual owned
elsewhere. This is the essential case: a borrowed endpoint may move across
processes because its protocol is a finite, bounded
obligation. Mobility extends from sessions to ordinary values: $\BUnit$ is mobile;
products and variants are mobile pointwise; a function value is mobile if its
arrow carries the mobile annotation, which in turn means the function 
closes over mobile values only. Context mobility is pointwise, so a process may
fork or send only closures whose captured assumptions are mobile.

A companion predicate $\Unr\T$ (and $\Unr\Ctx$ on contexts) identifies the
duplicable fragment: unit, componentwise-unrestricted products and variants,
and functions of direction $\DirNone$, lifted pointwise to contexts. It governs
contraction, weakening, and recursive-function introduction, and every
unrestricted context is also mobile. The full definition is shown in
\cref{fig:omitted-predicates} in the supplement.

\begin{figure}
{
  Equality of typing contexts is the equivalence closure of the following
  axioms.
  \smallskip
  \begin{mathpar}
    \CtxEmpty\CtxPar\Ctx = \Ctx \and
    \Ctx[1] \CtxPar \Ctx[2] = \Ctx[2] \CtxPar \Ctx[1] \and
    \Ctx[1] \CtxSeq (\Ctx[2] \CtxSeq \Ctx[3]) = (\Ctx[1] \CtxSeq \Ctx[2]) \CtxSeq \Ctx[3] \and
    \Ctx[1] \CtxPar (\Ctx[2] \CtxPar \Ctx[3]) = (\Ctx[1] \CtxPar \Ctx[2]) \CtxPar \Ctx[3] \\
    \inferrule{\Unr\Ctx}{\Ctx \CtxPar \Ctx = \Ctx} \and
    \inferrule{\Mobile\Ctx[1]}{\Ctx[1] \CtxPar \Ctx[2]
      = \Ctx[1] \CtxSeq \Ctx[2]} \and
    \inferrule{\Mobile\Ctx[2]}{\Ctx[1] \CtxPar \Ctx[2]
      = \Ctx[1] \CtxSeq \Ctx[2]}
  \end{mathpar}
}

  \bigskip
{
  \setlength{\belowdisplayskip}{0pt}
  The subcontext relation for weakening is the reflexive, transitive congruence
  closure\\ of these two axioms.
  \begin{mathpar}
    (\Ctx[1] \CtxPar \Ctx[3]) \CtxSeq (\Ctx[2] \CtxPar \Ctx[4])
    \SubCtx
    (\Ctx[1] \CtxSeq \Ctx[2]) \CtxPar (\Ctx[3] \CtxSeq \Ctx[4])

    \inferrule{
      \Unr\Ctx
    }{
      \CtxEmpty
      \SubCtx
      \Ctx
    }
  \end{mathpar}
}
  \caption{Equality of contexts and the subcontext relation $\Ctx\SubCtx\Ctx$}
  \label{fig:syntax-types}
\end{figure}

\Cref{fig:syntax-types} makes context equality sensitive to resource
discipline. Parallel composition is commutative and associative with empty
context as unit. Sequential composition is associative and preserves order.
Unrestricted contexts contract under $\CtxPar$, which is the only source of
implicit duplication. Mobile contexts may move between $\CtxPar$ and
$\CtxSeq$ connectives, reflecting that the type system is not
responsible for ordering access to mobile values.

The exchange axiom in the subcontext
relation  $\Ctx[1]\SubCtx\Ctx[2]$ is the concurrent-Kleene-algebra law that embeds two independent
ordered computations into one ordered computation over parallel
components \cite{DBLP:journals/jlp/HoareMSW11}.
The second axiom admits unused unrestricted assumptions. By
\RuleName{T-Weaken}, a derivation obtained in some context can be used in
any larger context related by~$\SubCtx$.

Session types themselves are quotiented by a standard equational theory used by
type conversion: sequential composition is associative with $\SSkip$ as neutral
element, recursion unfolds equi-recursively, and a continuation distributes
homomorphically through internal and external choice. The induced congruence is
the equivalence $\SEq$ that appears throughout the type system.

A fresh
channel offers dual protocols to its two endpoints. Duality, denoted $\SDual\T$, is the standard
involution that swaps send with receive and internal with external choice. We leave duality undefined on the synchronization markers $\SDrop$ and $\SAcq$, which never
occur on a $\SNew$ channel. Both the equational laws and the duality laws are
collected in \cref{fig:session-laws} in the supplement.

\begin{figure}
  \begin{mathpar}
    \RuleCTUnit \and
    \RuleCTFork \and
    \RuleCTNew \and
    \RuleCTLSplit \and
    \RuleCTRSplit \and
    \RuleCTDrop \and
    \RuleCTAcquire \and
    \RuleCTTerm \and
    \RuleCTWait \and
    \RuleCTSend \and
    \RuleCTRecv \and
    \RuleCTSelect \and
    \RuleCTBranch \and
    \RuleCTDiscard
  \end{mathpar}
  \caption{Constant typing ($\HasCType\C\T$)}
  \label{fig:constant-typing}
\end{figure}

\Cref{fig:constant-typing} assigns closed types to constants. Channel creation
produces a pair of dual endpoints, each protected by an initial $\SAcq$ marker
and closed respectively by $\STerm$ and $\SWait$. The session supplied to
$\CNew$ must be a source protocol: the side condition $\SNew\S$ restricts it to
skips, sends, receives, choices, branches, sequencing, and recursion, excluding
the termination markers $\STerm$, $\SWait$, and the synchronization markers
$\SAcq$, $\SDrop$, which are introduced only by remote splitting. Communication
constants consume exactly the session prefix they implement: send and receive
advance through $\SSend\T$ and $\SRecv\T$, selection advances an internal
choice, and branching returns a variant over external branches. These constants
are impure because they communicate.
%
The split constants are pure as they do not incur any communication by
themselves. Local splitting, $\CSplit$, decomposes $\SSeq{\S[1]}{\S[2]}$ into
an ordered pair containing the borrowed prefix and the residual suffix. The
ordering enforces that the prefix is consumed before communication operations
on the suffix become possible. Remote splitting, $\CSSplit$, produces an
unordered pair in which the prefix ends in $\SDrop$ and the residual begins
with $\SAcq$.
%
Both splits require the suffix $\S[2]$ to be nontrivial and local,
$\SLocal\S[2]$: it must not begin with $\SAcq$, looking through leading skips
and recursion binders. A formal definition is given in
\cref{fig:omitted-predicates}. In addition, $\CSplit$ does not allow a trivial
prefix either.
%
The constants $\CDrop$ and $\CAcq$ eliminate the
corresponding marker types. Forking expects a mobile impure computation from
unit to unit, and $\CDiscard$ consumes a trivial session.

% Example should be placed after the discussion of constants and their typing.
Consider $\FnRenderUser$ from \cref{lst:rendering-example}. Its full type is
$ \typecolor{\TyUser} \TArrx[m,u,p] \SSeq \TyHtmlChan \STerm
\TArrx[m,u,i] \BUnit $.  The first arrow is mobile ($\MobMobile$) and
unrestricted ($\DirNone$) because $\FnRenderUser$ is a top-level function,
which can be shared and duplicated at will. The arrow is also marked pure
($\EffPure$) because application of the $\TyUser$ argument returns a closure
without performing any channel communication. The closure contains only
unrestricted, mobile types which makes it unrestricted and mobile itself and is
reflected in the $\MobMobile$ and $\DirNone$ annotations on the second arrow.
The arrow is impure because application of the channel argument evaluates the
function body that contains calls to impure functions.

The difference in premises between $\CSplit$ and $\CSSplit$ enables the
latter---in conjunction with $\CDrop$---to transform any local, bounded channel
into a mobile one. A local channel $\EVarc$ of type $\S$ cannot cross thread
boundaries. The expression $\EApp {\CSSplit_{\SSeq\SSkip\S}} \EVarc$ gives a
pair of type $\EPair {\SDrop} {\SSeq\SAcq\S}$. The first component can be
dropped immediately; the second is the channel $\EVarc$ but mobile. 

\begin{figure}
  \begin{mathpar}
    \RuleTConst \and
    \RuleTVar \\
    \RuleTAbsUnr \and
    \RuleTAppUnr \and
    \RuleTAbsLin \and
    \RuleTAppLin \and
    \RuleTAbsLeft \and
    \RuleTAppLeft \and
    \RuleTAbsRight \and
    \RuleTAppRight \and
    \RuleTPairUnord \and
    \RuleTPairOrd \and
    \RuleTSeq \and
    \RuleTLetUnord \and
    \RuleTLetOrd \and
    \RuleTInj \and
    \RuleTCase \and
    \RuleTWeaken \and
    \RuleTConv
  \end{mathpar}
  \caption{Expression typing ($\HasType[\Eff]\Ctx\E\T$)}
  \label{fig:typing-declarative}
\end{figure}

The expression typing judgment $\HasType[\Eff]\Ctx\E\T$ in
\cref{fig:typing-declarative} records both resource usage and effects. The
context $\Ctx$ describes the assumptions consumed by $\E$ to compute a value of
type $\T$. The effect
annotation $\Eff$ is an upper bound, allowing pure expressions where impure
ones are expected.

Constants and variables are pure and use exactly the resources shown in their
types. Function introduction records how the argument is positioned relative
to the captured context. An unordered function places the argument in parallel
with the captured context, left- and right-ordered functions place it before or
after the captured context respectively and unrestricted recursive functions use direction
$\DirNone$ and may only capture an unrestricted context.
%
Application composes the contexts of the function and argument according to the
arrow direction. If the function is either $\DirLeft$- or $\DirRight$-ordered,
the function term (or correspondingly the argument term) must be pure.

Pair introduction mirrors function application. Unordered pairs combine their
components in parallel. Ordered pairs use a sequential context and require the
second component to be pure. Sequencing requires the subterm to be
unrestricted, so that its result can be ignored. Let and pair elimination expose bindings
to the body using the same ordered or unordered structure as the eliminated
type. Case analysis requires all branches to produce the same result type
under the same surrounding context shape. The final rules close the judgment
under weakening, type conversion, and effect subsumption.

The need for the purity constraints in the application of ordered functions and
the construction of ordered pairs arises from the difference between variable
usage and consumption of the actual resources. A channel is consumed, for
example, by passing it to $\CRecv$. One the other hand, applying the identity
function $\termcolor{\Varid{id}}$ to a variable already counts this as a usage
of said variable. The problem is demonstrated by the expression
%
\(
  \ELetPair[\DirLeft]{c_1}{c_2}{\EApp{\CSplit_{\SSeq{\S[1]}{\S[2]}}}{c}}
  \!\termcolor{\EPairOrd {\EApp{\Varid{id}}{c_1}} {\EApp{f}{c_2}}}
\).
%
The effect system ensures that $\EVarf$ does not communicate on $\EVarc[2]$.

% The function $\termcolor{f}$ must not be perform communication operations on
% $\termcolor{c_2}$ before the communication according to $\S[1]$ has taken place
% on $\termcolor{c_1}$.

\begin{figure}
  \begin{mathpar}
    \RuleTExpr \and
    \RuleTPar \and
    \RuleTRes \\
    \RuleBindEmp \and
    \RuleBindSeq \and
    \RuleBindDrop \and
    \RuleBindAcq
  \end{mathpar}
  \caption{Process typing ($\ProcHasType \Ctx {\Proc}$) and binding
    ($\BindGivesCtx {\S} {\BindGroup} {\Ctx}$)}
  \label{fig:process-typing}
\end{figure}

Processes are typed by $\ProcHasType\Ctx\Proc$ in
\cref{fig:process-typing}. An expression process may be an impure computation
of type unit. Parallel composition splits the process context with $\CtxPar$.
Restriction introduces two binder groups at dual session types and adds
the assumptions generated for those binders in parallel to the subprocess.

The auxiliary judgment $\BindGivesCtx\S\BindGroup\Ctx$ relates binder
groups to contexts. A plain sequence of variables follows the sequential
structure of the session type. A separator splits a session
$\SSeq{\S[1]}{\S[2]}$ into a prefix ending in $\SDrop$ and a suffix guarded by
$\SAcq$, and places the resulting contexts in parallel. Thus binder groups are
the static representation of the synchronization boundary introduced by
remote borrowing. This judgment is nondeterministic as session types
are subject to conversion.

%%% Local Variables:
%%% mode: latex
%%% TeX-master: "../popl2027.tex"
%%% End:
